Describe the setting of the Generalized Linear Model (GLM)

\( g(\mathbb{E}[y|X]) = X\beta = \nu\)

We will denote as usual \( \mu = \mathbb{E}[y|X] \).
\( \nu \) is called the linear predictor.
The choice of \( g(\mu) \), called the link, determines how the mean is related
to the explanatory variables \( x \). \( g^{-1} \) is known as activation function in ML.

In the classical linear model, \( g = \text{id} \).
In the GLM, this is generalized to a \( g \) which is a monotonic,
differentiable function (such as log or square root).

Given \( X \), \( \mu \) is determined through \( g(\mu) \). 
Given \( \mu, \theta \) is determined through \( \dot{a}(\theta) = \mu \). 
Finally given \( \theta \), the outcomes \( y \) are determined as a draw 
from the exponential density specified in \( a(\theta) \).

Give a list of link functions and explain offsets

INSERT IMAGE FROM 67


Modeling counts such as the number of claims or deaths in a risk group
requires correction for the number \( n \) exposed to risk. If \( \mu \) is the mean of the
count \( y \), then the occurrence rate \( \frac{\mu}{n} \) of interest and \( g(\frac{\mu}{n}) = X\beta \).
When \( g = \log \) then \( \log(\frac{\mu}{n}) = X\beta \equiv \log(\mu) = \log(n) + X\beta \).


Derive the MLE equations for the GLM

\(\ell(\beta, \phi) = \sum_{i=1}^{n} \ln f(y_i; \beta, \phi) = \sum_{i=1}^{n} \left\{ \ln c(y_i, \phi) + \frac{y_i \theta_i - a(\theta_i)}{\phi} \right\},\)

which assumes independent exponential family responses \(y_i\). Consider the MLE of \(\beta_j\). To find the maximum, \(\ell(\beta, \phi)\) is differentiated with respect to \(\beta_j\): 
\(\frac{\partial \ell}{\partial \beta_j} = \sum_{i=1}^{n} \frac{\partial \ell}{\partial \theta_i} \frac{\partial \theta_i}{\partial \beta_j}.\)
where
\[
\frac{\partial \ell}{\partial \theta_i} = \frac{y_i - \dot{a}(\theta_i)}{\phi} = \frac{y_i - \mu_i}{\phi}, \quad \frac{\partial \theta_i}{\partial \beta_j} = \frac{\partial \theta_i}{\partial \eta_i} \frac{\partial \eta_i}{\partial \beta_j} = \frac{\partial \theta_i}{\partial \eta_i} x_{ij}.
\]

Here \(\eta_i = x_i' \beta\) and \(x_{ij}\) is component \(i\) of \(x_j\). Setting \(\partial \ell/\partial \beta_j = 0\) yields 
the first order conditions for likelihood maximization:
\[
\sum_{i=1}^{n} \frac{\partial \theta_i}{\partial \eta_i} x_{ij} (y_i - \mu_i) = 0 \quad \Leftrightarrow \quad X'D(y - \mu) = 0,
\]
where \(D\) is the diagonal matrix with diagonal entries \(\partial \theta_i/\partial \eta_i\),
\[
\left(\frac{\partial \theta_i}{\partial \eta_i}\right)^{-1} = \frac{\partial \eta_i}{\partial \theta_i} = \frac{\partial \eta_i}{\partial \mu_i} \frac{\partial \mu_i}{\partial \theta_i} = g'(\mu_i) \dot{a}(\theta_i) = g'(\mu_i) V(\mu_i).
\]
Thus \(D\) is diagonal with entries \(\{g'(\mu_i) V(\mu_i)\}^{-1}\). 
The equations in (5.4) are often called the estimation equations for \(\beta\). 
Note \(\beta\) is implicit in these equations, working implicitly through \(\mu\) and \(D\).

Defining the diagonal matrices \(G\) and \(W\) with diagonal entries \(g'(\mu_i)\) and \(\{g'(\mu_i)^2 V(\mu_i)\}^{-1}\), respectively, 
then \(D = WG\) and (5.4) is equivalent to
\[
X'WG(y - \mu) = 0.
\]


Describe the Newton-Raphson method for likelihood maximization

Suppose \(\phi\) is known and so write \(\ell(\beta)\) to denote the log-likelihood as a function of the unknown parameter vector \(\beta\). 
If \(\beta\) contains a single parameter then the quadratic Taylor series approximation at any point \(\beta\) is 
\(\ell(\beta + \delta) \approx \ell(\beta) + \dot{\ell}(\beta)\delta + \frac{\delta^2}{2} \ddot{\ell}(\beta).\)

Differentiating the right-hand side as a function of \(\delta\) and equating to zero yields 
\(\dot{\ell}(\beta) + \delta \ddot{\ell}(\beta) = 0 \quad \Rightarrow \quad \delta = - \{\ddot{\ell}(\beta)\}^{-1} \dot{\ell}(\beta).\)

With \(\beta\) given and \(\delta\) as specified, a higher point thus appears to be \(\beta - \dot{\ell}(\beta)/\ddot{\ell}(\beta)\). 
Denoting \(\beta^{(m)}\) as the value for \(\beta\) at iteration \(m\), the update equation is
\(\beta^{(m+1)} = \beta^{(m)} - \{\ddot{\ell}(\beta^{(m)})\}^{-1} \dot{\ell}(\beta^{(m)}). \qquad (5.8)\)

For a maximum \(\ddot{\ell}(\beta) < 0\). 
Iteration of this equation, each time revising \(\beta\), leads to a sequence which, as stated above, often rapidly converges.

The approach can be adapted when \(\beta\) is vector. In this case a quadratic approximation is made to the surface \(\ell(\beta)\), 
and it is this quadratic surface that is maximized. The update equation is as in (5.8) with the inverse interpreted as matrix inversion 
and where \(\dot{\ell}(\beta)\) is the vector of partial derivatives \(\partial \ell / \partial \beta_k\). 

The vector \(\dot{\ell}(\beta)\) is called the "score" vector, and \(\ddot{\ell}(\beta)\), the matrix of cross partial derivatives 
\(\partial^2 \ell/\{\partial \beta_j \partial \beta_k\}\), the "Hessian" matrix. 
The condition for a maximum is that the Hessian is non-positive definite, i.e., \(-\ddot{\ell}(\beta)\) is non-negative definite. 
The procedure of repeated application of the score and Hessian matrix to equation (5.8) is called score function iteration.
